\section{The loop on digital flow}

\begin{frame}{Greedy loop, applied to the digital flow}
  \begin{center}
  \begin{tikzpicture}[node distance=0.5cm and 0.7cm]
    \node[agent3, line width=0.8mm] (prop) {Propose\\\scriptsize LLM (cc\_cli or opencode)};
    \node[agent, right=of prop] (map) {Map params\\$\rightarrow$ \cmd{config.yaml}};
    \node[agent, right=of map] (eval) {Evaluate\\\scriptsize LibreLane + cocotb GL};
    \node[agent, below=of eval] (fom) {Score FoM\\\scriptsize weighted, valid-gated};
    \node[agent, left=of fom] (log) {Log to\\\cmd{results.tsv}};
    \node[agent, left=of log] (prog) {Append to\\\cmd{program.md}};
    \draw[arrow] (prop) -- (map);
    \draw[arrow] (map) -- (eval);
    \draw[arrow] (eval) -- (fom);
    \draw[arrow] (fom) -- (log);
    \draw[arrow] (log) -- (prog);
    \draw[arrow] (prog) |- (prop);
  \end{tikzpicture}
  \end{center}
  \vspace{0.3em}
  \footnotesize
  \begin{itemize}
    \item Same Karpathy-autoresearch shape as analog; only the
          \term{evaluator} and \term{design space} change.
    \item The \warn{Propose} node is what dispatches by
          \cmd{--backend} (covered in section 4).
    \item Strategy axis (commit \cmd{7b5720d}):
          \cmd{flow} (knobs only), \cmd{rtl} (RTL micro-edits),
          \cmd{hybrid} (both). This deck focuses on \cmd{flow}.
  \end{itemize}
  \note{%
    Reuse of the analog loop figure is intentional: same dispatcher,
    same persistence pattern, same budget knob. The two things that
    change at the digital plug-in: the evaluator runs a full bare
    LibreLane flow plus post-synth and post-PnR-SDF cocotb GL sim
    instead of ngspice; the design space is discrete sets of flow
    knobs rather than continuous W/L. \cmd{program.md} stays the
    LLM's narrative memory across iterations so the prompt context
    stays small.}
\end{frame}

\begin{frame}{The same loop, framed as emulated RL}
  \begin{center}
  \begin{tikzpicture}[every node/.style={font=\footnotesize}]
    \node[notebox, minimum width=7.4cm, minimum height=0.9cm]
      (mem) at (0, 2.0)
      {\textbf{Memory} (in-context experience replay)\\
       \cmd{program.md} \,+\, \cmd{results.tsv}};
    \node[agent3, minimum width=3.2cm, minimum height=1.4cm,
          line width=0.8mm] (ag) at (-3.7, 0)
      {\textbf{Agent / policy}\\
       \scriptsize pretrained LLM\\
       \scriptsize \warn{frozen weights}};
    \node[agent, minimum width=3.2cm, minimum height=1.4cm]
      (env) at (3.7, 0)
      {\textbf{Environment}\\
       \scriptsize LibreLane flow\\
       \scriptsize + cocotb GL sim};
    \draw[arrow] (mem.south west) to[bend right=20]
      node[midway, left=1pt, font=\scriptsize, text=muted, align=center]
      {context\\replay} (ag.north);
    \draw[arrow] (env.north) to[bend right=20]
      node[midway, right=1pt, font=\scriptsize, text=muted, align=center]
      {append\\eval row} (mem.south east);
    \draw[arrowclaude] (ag.east) to[bend left=15]
      node[midway, above=1pt, font=\scriptsize]
      {action: \cmd{params} dict} (env.west);
    \draw[arrow2] (env.west) to[bend left=15]
      node[midway, below=1pt, font=\scriptsize]
      {observation + reward (FoM)} (ag.east);
  \end{tikzpicture}
  \end{center}
  \vspace{-0.2em}
  \footnotesize
  \begin{itemize}
    \item \term{Action}: flow-knob \cmd{params} dict.
          \term{Reward}: scalar FoM, gated to $0$ on any signoff
          failure.
    \item \warn{Why emulated}: no gradient step; in-context
          replay (\cmd{program.md} + last rows of
          \cmd{results.tsv}) substitutes for weight updates.
          Episode cap = \cmd{--budget}.
  \end{itemize}
  \note{%
    The architecture matches RL exactly -- agent, environment,
    reward, experience -- but the learning algorithm is replaced
    by in-context retrieval. Real RL is impractical here because
    each environment step costs 8-25 minutes of LibreLane wall
    time, so collecting the millions of samples a DQN or PPO
    needs is out of reach. Pretrained LLMs already encode strong
    priors over EDA-flow trade-offs (clock vs density vs synth
    strategy), so a handful of evals usually suffices. This is
    also why \cmd{program.md} matters: it is the LLM's
    hand-written replay buffer, deliberately kept small enough to
    fit comfortably in the prompt context. \cmd{results.tsv}
    carries the numeric history; the LLM rarely reads more than
    the last few rows.}
\end{frame}

\begin{frame}{What gets proposed, what gets scored on GF180}
  \driving{LibreLane + cocotb GL sim}
  \begin{columns}[T,onlytextwidth]
    \column{0.48\textwidth}
    \textbf{\color{accent}Proposed (flow strategy)}
    \footnotesize
    \begin{itemize}
      \item \cmd{CLOCK\_PERIOD}
      \item \cmd{PL\_TARGET\_DENSITY\_PCT}
      \item \cmd{SYNTH\_STRATEGY}
      \item PDN tweaks (rail width, ring count)
    \end{itemize}
    \vspace{0.3em}
    \scriptsize
    \cmd{rtl} and \cmd{hybrid} strategies also propose
    free-form RTL micro-edits (resource sharing, FSM re-encoding).

    \column{0.5\textwidth}
    \textbf{\color{accent2}Scored (FoM)}
    \footnotesize
    \[
      \mathrm{FoM} =
        w_t\,s_{\text{timing}} + w_p\,s_{\text{perf}}
      + w_a\,s_{\text{area}}  + w_e\,s_{\text{energy}}
    \]
    \scriptsize
    \begin{itemize}
      \item \cmd{flow\_metrics.py:213}; default $w_t=0.5$, others~$1.0$.
      \item $s_{\text{timing}}$ is binary; \warn{any DRC/LVS/sim
            failure sets} \cmd{valid=False} \warn{and FoM=0}.
      \item No \cmd{--stop-after}: signoff stays on every eval.
    \end{itemize}
  \end{columns}
  \note{%
    The valid gate is the unsugared part: a proposal is only
    \emph{kept} when LibreLane closes timing, DRC and LVS report
    zero violations, and both post-synth and post-PnR-SDF cocotb
    GL sim runs pass against the design's testbench. This is why
    \cmd{DigitalDesign.testbench()} became mandatory (commit
    \cmd{e386cf2}); without it the loop has no gate to fail
    against. We deliberately keep no \cmd{--stop-after} knob; the
    point of the loop is to optimise on real signoff metrics, not
    on intermediate proxies. The downside is wall time: 8-25 min
    per eval on GF180 for a small block.}
\end{frame}
